\documentclass[11pt,a4paper]{article}
%% Shared SCRS preamble fragments (input after \documentclass)
\usepackage[margin=2.5cm]{geometry}
\usepackage{setspace}
\usepackage{authblk}
\usepackage{graphicx}
\usepackage{booktabs}
\usepackage{array}
\usepackage{placeins}
\usepackage{float}
\usepackage{amsmath,amssymb}
\usepackage{enumitem}
\usepackage{url}
\usepackage[round]{natbib}
\usepackage{hyperref}
\hypersetup{colorlinks=true,linkcolor=black,citecolor=black,urlcolor=blue}

\setstretch{1.0}
\parskip=6pt
\parindent=0pt
\pagestyle{empty}
\emergencystretch=3em
\graphicspath{{figs/}}

\newcommand{\incfig}[2][]{%
  \IfFileExists{figs/#2}{%
    \includegraphics[#1]{#2}%
  }{%
    \fbox{\begin{minipage}{0.92\linewidth}\centering\vspace{1.5em}
      Missing figure: \texttt{#2}\\[0.5em]
      {\small regenerate from \texttt{Rmd/sma-uncertainty.Rmd}}
      \vspace{1.5em}\end{minipage}}%
  }%
}

% generated by scripts/refresh-paper-numbers.R; do not edit
\newcommand{\nDraw}{1000}
\newcommand{\jabbaRun}{JABBA3}
\newcommand{\histEnd}{2024}
\newcommand{\linN}{20{,}000}
\newcommand{\linCells}{60}
\newcommand{\linMaxDiff}{4}
\newcommand{\linMeanDiff}{1.0}
\newcommand{\linWithinOne}{45}
\newcommand{\linMaxRounded}{4}
\newcommand{\trDiff}{3.3}
\newcommand{\statusGreenMvln}{52}
\newcommand{\statusGreenMcmc}{59}
\newcommand{\statusGreenPe}{79}
\newcommand{\statusGreenJabba}{76}
\newcommand{\statusGreenJabbaDyn}{86}
\newcommand{\histBMvln}{1.01 (0.80--1.30)}
\newcommand{\histFMvln}{0.73 (0.51--1.08)}
\newcommand{\histBMcmc}{1.11 (0.84--1.46)}
\newcommand{\histFMcmc}{0.96 (0.67--1.35)}
\newcommand{\histBMvlnMc}{n/a}
\newcommand{\histFMvlnMc}{n/a}
\newcommand{\histBJabba}{1.22 (0.67--2.18)}
\newcommand{\histFJabba}{0.29 (0.07--1.19)}
\newcommand{\histBJabbaDyn}{1.31 (0.82--2.14)}
\newcommand{\histFJabbaDyn}{0.22 (0.05--0.87)}
\newcommand{\histBJabbaOneB}{1.31 (0.82--2.14)}
\newcommand{\histFJabbaOneB}{0.22 (0.05--0.87)}
\newcommand{\histBJabbaOneS}{1.22 (0.67--2.18)}
\newcommand{\histFJabbaOneS}{0.29 (0.07--1.19)}
\newcommand{\histBJabbaTwoS}{n/a}
\newcommand{\histFJabbaTwoS}{n/a}
\newcommand{\greenZeroMvlnFortyFive}{96}
\newcommand{\tacMvlnFortyFiveSixty}{1,288}
\newcommand{\tacMvlnFortyFiveSeventy}{1,033}
\newcommand{\greenZeroMvlnSeventy}{100}
\newcommand{\tacMvlnSeventySixty}{1,587}
\newcommand{\tacMvlnSeventySeventy}{1,486}
\newcommand{\greenZeroMcmcFortyFive}{99}
\newcommand{\tacMcmcFortyFiveSixty}{1,185}
\newcommand{\tacMcmcFortyFiveSeventy}{1,080}
\newcommand{\greenZeroMcmcSeventy}{100}
\newcommand{\tacMcmcSeventySixty}{1,188}
\newcommand{\tacMcmcSeventySeventy}{1,114}
\newcommand{\greenZeroPeFortyFive}{100}
\newcommand{\tacPeFortyFiveSixty}{1,320}
\newcommand{\tacPeFortyFiveSeventy}{1,150}
\newcommand{\greenZeroPeSeventy}{100}
\newcommand{\tacPeSeventySixty}{1,300}
\newcommand{\tacPeSeventySeventy}{1,217}
\newcommand{\greenZeroJabbaFortyFive}{97}
\newcommand{\tacJabbaFortyFiveSixty}{2,300}
\newcommand{\tacJabbaFortyFiveSeventy}{2,300}
\newcommand{\greenZeroJabbaSeventy}{98}
\newcommand{\tacJabbaSeventySixty}{2,300}
\newcommand{\tacJabbaSeventySeventy}{2,300}
\newcommand{\greenZeroJabbaDynFortyFive}{98}
\newcommand{\tacJabbaDynFortyFiveSixty}{2,300}
\newcommand{\tacJabbaDynFortyFiveSeventy}{2,300}
\newcommand{\greenZeroJabbaDynSeventy}{100}
\newcommand{\tacJabbaDynSeventySixty}{2,300}
\newcommand{\tacJabbaDynSeventySeventy}{2,300}
\newcommand{\nMcmc}{1,485}
\newcommand{\mcmcCatches}{26}
\newcommand{\apxMeanMvln}{n/a}
\newcommand{\apxMaxMvln}{n/a}
\newcommand{\apxMeanMcmc}{n/a}
\newcommand{\apxMaxMcmc}{n/a}
\newcommand{\apxMeanJabba}{n/a}
\newcommand{\apxMaxJabba}{n/a}
\newcommand{\tacRangeSixtySeventy}{713}
\newcommand{\rpBaseSsbmsy}{552}
\newcommand{\rpBaseMsy}{1,880}
\newcommand{\rpMleSsbmsy}{360}
\newcommand{\rpMleMsy}{1,240}
\newcommand{\rpMcSsbmsy}{369}
\newcommand{\rpMcMsy}{1,280}
\newcommand{\rpSsbmsyDiff}{33}
\newcommand{\rpMsyDiff}{32}
\newcommand{\rpFDiff}{29}

\InputIfFileExists{numbers-ssproject.tex}{}{}

\newcommand{\PG}{P(\mathrm{green})}
\newcommand{\BB}{B/B_{\mathrm{MSY}}}
\newcommand{\SB}{\mathrm{SSF}/\mathrm{SSF_{MSY}}}
\newcommand{\FF}{F/F_{\mathrm{MSY}}}

\date{}

\renewcommand\thesection{\arabic{section}.}
\renewcommand\thesubsection{\arabic{section}.\arabic{subsection}}
\makeatletter
\renewcommand\section{\@startsection{section}{1}{\z@}%
  {-1.5ex \@plus -1ex \@minus -.2ex}{0.8ex \@plus .2ex}{\normalfont\bfseries}}
\renewcommand\subsection{\@startsection{subsection}{2}{\z@}%
  {-1.25ex \@plus -1ex \@minus -.2ex}{0.6ex \@plus .2ex}{\normalfont\itshape}}
\makeatother


\title{\textbf{Kobe Strategy Matrix Projections: Comparing\\
SS3-MVLN, SS3-MCMC, JABBA and JABBA with state-dependent process error\\
\large (worked example: North Atlantic shortfin mako)}\\[0.5em]
\large SCRS/2026/XXX --- for WGSAM / Shark WG}

\author[1]{L. T. Kell}
\affil[1]{Centre for Environmental Policy, Imperial College London, London, United Kingdom.\\
\texttt{laurie@seaplusplus.co.uk}}

\begin{document}
\maketitle

\begin{center}\textit{SUMMARY}\end{center}

This paper records the four constant-catch projections used to apply
rebuilding-time reference points to North Atlantic shortfin mako (SMA-N).
The stock is a non-retained bycatch under Rec.~21-09, which requires
rebuilding with probability 60--70\% by 2070. The same catch grid, years and
definition of Kobe green are used for SS3-MVLN, SS3-MCMC, JABBA3 with
constant process error, and JABBA3 with state-dependent process error. Green
is $B \ge B_{\mathrm{MSY}}$ and $F \le F_{\mathrm{MSY}}$, with biomass at the
end of the year. For SS3-MVLN that probability is the bivariate normal
probability written in Section~\ref{sec:mvln}; the numbers below are a
sample of \nDraw{} draws from that distribution. The other three
projections count their own draws. The reported products
are the catch at $\PG = 0.6$ and $0.7$ in 2045 and 2070, and the rebuilding
year at each catch. On this grid the two JABBA3 projections remain above both
targets at 2{,}300~t, so their risk-equivalent catch lies above the grid,
while the SS3-MVLN catch at 60\% in 2070 is \tacMvlnSeventySixty~t. The
SS3-MCMC contrast still uses a different recruitment constraint
(\texttt{do\_recdev} = 2) and is illustrative only. The JABBA3 runs are an
SS3-to-surplus-production bridge (\texttt{ss2jabba3}); they are not the Shark
WG JAGS scenarios. The paper does not set Species Group catch advice.

\textbf{KEYWORDS:} rebuilding, Kobe strategy matrix, MVLN, MCMC, JABBA3,
process error, shortfin mako, non-retained bycatch, Rec.~21-09

\section{Introduction}

Rebuilding-time reference points fix a probability and a year, then read the
catch, or fix the catch and read the year \citep{KellAcceptedFaF}. That is
the forward-looking form of risk equivalence \citep{Fischer2023}. For SMA-N
the probability and the year are already set: Rec.~21-09 requires $\PG = 0.6$
to $0.7$ by 2070 for a stock whose retention is prohibited
\citep{ICCAT2021Rec2109,ICCAT2022Rec2211}. The object that has to be inverted
is a Kobe strategy matrix.

The 2026 assessment provides the SS3 base case and a Kobe matrix filled by
drawing from the multivariate log-normal \citep{Liniers2026}. This paper
places four constant-catch projections on one grid, so the catch at
$\PG = 0.6$ and $0.7$, and the year at which each catch meets those
probabilities, can be compared. The four are SS3-MVLN, SS3-MCMC, JABBA3
with constant process error, and JABBA3 with state-dependent process error.
The paper specifies how each projection fills
$\PG(\mathrm{catch}, \mathrm{year})$ and reports the catches and rebuilding
years. It does not set Species Group catch advice. The endorsed projection
remains the SS3 run adopted by the Shark WG.

\section{Shared design}

One stock, the 2026 SS3 base configuration
\texttt{HW2e\_HW2d\_InitFprior} \citep{Liniers2026}. One constant-catch grid:
0 to 2{,}300~t in steps of 100~t, plus 250~t and 1{,}883~t. Projection years
2025--2070; the assessment ends in \histEnd. Two target probabilities, 0.6
and 0.7. Two years at which the catch is read, 2045 and 2070.

Green is the joint event $B \ge B_{\mathrm{MSY}}$ and
$F \le F_{\mathrm{MSY}}$. For SS3, $B$ is end-of-year spawning stock fecundity
(SSF), paired with $F$ of the same year, as in the 2026 Kobe matrix. For
JABBA3, $B$ is surplus-production biomass. Absolute tonnes need not match
across those scales. The comparison uses $\BB$, $\FF$ and $\PG$.

For each projection and each catch, $\PG$ is interpolated on catch at the
target year. The catch at which it equals 0.6 or 0.7 is the risk-equivalent
catch. If $\PG$ is still above the target at 2{,}300~t, the table records
2{,}300~t and the text states that the crossing lies above the grid. The
rebuilding year is the last time $\PG$ crosses the target, provided the
series ends at or above the target (\texttt{kobeRebuildTime}). A first
crossing is not used: for a late-maturing stock $\PG$ can dip and recover,
and an early crossing would mark a stock that is still on the way down.

\begin{table}[htbp]
\caption{\label{tab:four}The four projections. They share the catch grid, the
years and the definition of green. They differ in the model and in how
$\PG$ is obtained.}
\centering
\begin{tabular}[t]{>{\raggedright\arraybackslash}p{0.22\textwidth}>{\raggedright\arraybackslash}p{0.34\textwidth}>{\raggedright\arraybackslash}p{0.34\textwidth}}
\toprule
Projection & Model & How $\PG$ is obtained\\
\midrule
SS3-MVLN & SS3, \texttt{do\_recdev}${}=1$ & Closed-form bivariate normal\\
SS3-MCMC & SS3, \texttt{do\_recdev}${}=2$ & Proportion of posterior draws\\
JABBA3 & Pella--Tomlinson, constant $\sigma$ & Proportion of \texttt{jb3\_project} draws\\
JABBA3 dyn & Same production function, $\sigma_t$ state-dependent & Proportion of \texttt{jb3\_project} draws\\
\bottomrule
\end{tabular}
\end{table}

Figures and the numbers quoted below are written by
\texttt{Rmd/sma-uncertainty.Rmd}. The shared knit also draws the recDev2
MVLN (\texttt{do\_recdev} = 2) and an \texttt{FLasher} process-error
series. Those two series appear in the trajectory, matrix, risk and
rebuilding figures. They are not among the four projections in
Table~\ref{tab:four}, and they are not used in Table~\ref{tab:tacfour}.

\section{SS3-MVLN}
\label{sec:mvln}

The data-rich delta-method projection uses the recDev1 base case, in which
recruitment deviations are constrained to sum to zero (\texttt{do\_recdev}
= 1). \texttt{ss3diags2::ss3mvln} and \texttt{ss3stk} supply, for each catch
and year, the log-scale means $\mu_B$ and $\mu_F$, standard deviations
$\sigma_B$ and $\sigma_F$, and correlation $\rho$ of end-of-year SSF and $F$
\citep{ss3diags2,Kell2023MVLN}. With $\Phi$ the standard normal distribution
function and $\Phi_2$ the standard bivariate normal distribution function,
\begin{align}
  P(F \le F_{\mathrm{MSY}}) &= \Phi\!\left(-\mu_F/\sigma_F\right), \\
  P(B \ge B_{\mathrm{MSY}}) &= \Phi\!\left(\mu_B/\sigma_B\right), \\
  \PG &= \Phi_2\!\left(\mu_B/\sigma_B,\; -\mu_F/\sigma_F;\; -\rho\right).
\end{align}
The log-scale correlation used here is the exact transform
$\log(1 + \rho\,\mathrm{CV}_B\mathrm{CV}_F)$. These values are exact for the
multivariate log-normal, so the matrix does not have to be simulated.
The catches and the strategy matrix printed below are the knit's sample of
\nDraw{} draws from that distribution. Near one half, the standard error of
a proportion from \nDraw{} draws is about 1.6 percentage points.
Appendix~\ref{app:mvln} compares the closed form itself with the published
SMA-N matrix, which was filled by drawing from the same distribution.

\section{SS3-MCMC}

The posterior projection uses the same TAC grid, re-evaluated with
\texttt{ss -mceval}. $\PG$ is the proportion of saved draws in the green
quadrant. The chain used here has \nMcmc{} draws after thinning, across
\mcmcCatches{} catches. Within each draw, SSF, $F$ and catch are scaled by
that draw's own $\mathrm{SSF_{MSY}}$, $F_{\mathrm{MSY}}$ and MSY.

The available chain was run with unconstrained recruitment deviations
(\texttt{do\_recdev} = 2). The MVLN base case uses \texttt{do\_recdev} = 1.
Reference points therefore differ between the two SS3 fits
(Table~\ref{tab:refpts}): the posterior median of $\mathrm{SSF_{MSY}}$
differs from the base-case MLE by \rpSsbmsyDiff\%, and MSY by \rpMsyDiff\%.
The base-case MLE has $\mathrm{SSF_{MSY}} = \rpBaseSsbmsy$ and
MSY $= \rpBaseMsy$~t; the MCMC model is centred near
$\mathrm{SSF_{MSY}} = \rpMcSsbmsy$ and MSY $= \rpMcMsy$~t. Differences
between SS3-MVLN and SS3-MCMC on this grid confound the uncertainty method
with that constraint. They are reported so the projection can be inspected,
and they are not a same-model test of delta method against MCMC
\citep{SharkWG2026}.

\begin{table}[H]

\caption{\label{tab:refpts}SS3 reference points and status: maximum likelihood estimate (MLE) of the base case (\texttt{do\_recdev} = 1), used for MVLN; MLE of the model from which the MCMC was run (\texttt{do\_recdev} = 2); and SS3 MCMC posterior median and 90\% interval.}
\centering
\begin{tabular}[t]{lrrrrr}
\toprule
Quantity & Base MLE & MCMC-model MLE & Median & 5\% & 95\%\\
\midrule
SSB\_MSY & 552 & 360 & 369 & 329 & 435\\
MSY (t) & 1,880 & 1,240 & 1,280 & 1,150 & 1,500\\
F\_MSY & 0.0129 & 0.013 & 0.013 & 0.0126 & 0.0134\\
SSB\_0 & 1,270 & 830 & 851 & 756 & 1,000\\
F/FMSY 2024 & 0.735 & 0.996 & 0.945 & 0.695 & 1.22\\
\addlinespace
SSB/SSBMSY end-2024 & 1.01 & 1.09 & 1.11 & 0.885 & 1.4\\
\bottomrule
\end{tabular}
\end{table}


\section{JABBA3 with constant process error}

The surplus-production tier is a Pella--Tomlinson model in JABBA3
\citep{JABBA3,Winker2018}. \texttt{ss3diags2::ss2jabba3} reads catch in
weight and the CPUE series from the recDev1 SS3 report. The production shape
$m$ is taken from \texttt{ss3diags2::ss3\_spm}. SMA spawning output is in
numbers, so $r$ and $K$ on that curve are not used as biomass priors. The
$r$ prior is log-normal with median 0.08 and log-scale standard deviation
0.4. The shape prior is log-normal, centred on the SS3 production-curve $m$,
with log-scale standard deviation 0.15. Process error has constant standard
deviation (\texttt{sdb.form = "const"}). The fit uses \nDraw{} draws. Constant
catches are projected with \texttt{jb3\_project} (\texttt{imp.years} = 2) and
$\PG$ is the proportion of draws in the green quadrant
(\texttt{jb3\_kobematrix}).

This is a reproducible bridge from the SS3 observation extract. It is not a
re-run of the assessment WG's JAGS JABBA scenarios, and it is not a candidate
to replace the endorsed SS3 projection while recovery depends on the lag
between exploitable biomass and spawners \citep{Winker2020SMA,SharkWG2026}.

\section{JABBA3 with state-dependent process error}

The second JABBA3 projection keeps the same production function and the same
priors. Process error is state-dependent,
\begin{equation}
\sigma_t = \sigma_{\mathrm{MSY}}\left(B_t/B_{\mathrm{MSY}}\right)^{-k},
\end{equation}
from \texttt{jb3\_tier2} with no maturity lag, so the lag in the production
function is unchanged. $\sigma_{\mathrm{MSY}}$ is scaled so that the average
process variance matches the constant-SD fit, and $k$ follows the life-history
prior in \texttt{jb3\_tier2}. The projection uses the same catch grid, the
same \nDraw{} draws and the same two-year implementation lag. The only
designed difference from the constant-SD run is the state dependence of
$\sigma_t$.

\section{Results}

\subsection{Status at the start of the projection}
In \histEnd{} the median $\BB$ (95\% interval) was \histBMvln{} under
SS3-MVLN, \histBMcmc{} under SS3-MCMC, \histBJabba{} under JABBA3 and
\histBJabbaDyn{} under JABBA3 dyn. The corresponding $\FF$ were
\histFMvln, \histFMcmc, \histFJabba{} and \histFJabbaDyn.
Figure~\ref{fig:phase} is the joint distribution of those draws in \histEnd,
with the marginal densities of $\BB$ and $\FF$.
In the first projection year, at the lowest catch, $\PG$ was
\statusGreenMvln\% (SS3-MVLN), \statusGreenMcmc\% (SS3-MCMC),
\statusGreenJabba\% (JABBA3) and \statusGreenJabbaDyn\% (JABBA3 dyn).

\begin{figure}[H]
  \centering
  \incfig[width=\textwidth]{hist-fig-1.png}
  \caption{Historical $\BB$ (top) and $\FF$ (bottom): median and 95\%
    interval. The panel on the right of each row is the marginal
    distribution in \histEnd. SS3 biomass is end-of-year SSF. SS3-MVLN uses
    \texttt{do\_recdev} = 1; SS3-MCMC uses \texttt{do\_recdev} = 2.}
  \label{fig:hist}
\end{figure}

\begin{figure}[H]
  \centering
  \incfig[width=0.78\textwidth]{phase-mar-1.png}
  \caption{Joint distribution of $\BB$ and $\FF$ in \histEnd, with marginal
    densities. Each cloud is a sample of the historical draws for that
    projection. The vertical and horizontal lines are $B_{\mathrm{MSY}}$ and
    $F_{\mathrm{MSY}}$.}
  \label{fig:phase}
\end{figure}

\begin{figure}[H]
  \centering
  \incfig[width=\textwidth]{status-bar-1.png}
  \caption{Kobe quadrant probabilities in \histEnd{} for the four
    projections that have a historical draw. Each bar is the proportion of
    draws in that quadrant. This is assessment-year status. The
    first-year projection probabilities are the percentages in the text.}
  \label{fig:status}
\end{figure}

\subsection{Projected biomass and fishing mortality}
Figures~\ref{fig:projb} and~\ref{fig:projf} run from 2000 through the
projection. Columns are the constant catches 0, 250, 1{,}200, 1{,}883 and
2{,}300~t. Rows, from top to bottom, are SS3 recDev1 (the SS3-MVLN
projection), the recDev2 MVLN, SS3-MCMC, FLasher, JABBA3 and JABBA3 dyn.
Ribbons are the 50, 80 and 95\% intervals and the line is the median. The
dashed line is $B_{\mathrm{MSY}}$ or $F_{\mathrm{MSY}}$. The dotted vertical
line is the last assessment year. MVLN draws are independent by year, so
the two MVLN rows have no line joining one draw through time. MCMC, FLasher
and both JABBA3 rows join the same draw from the historical years into the
projection. FLasher history is the MCMC sample that conditions that
projection.

\begin{figure}[H]
  \centering
  \incfig[width=\textwidth]{proj-b-1.png}
  \caption{Historical and projected $\BB$ by series and constant catch.
    Layout and line conventions are given in the text. The catch table uses
    only the four projections in Table~\ref{tab:four}.}
  \label{fig:projb}
\end{figure}

\begin{figure}[H]
  \centering
  \incfig[width=\textwidth]{proj-f-1.png}
  \caption{Historical and projected $\FF$. Same rows, catches and intervals
    as Figure~\ref{fig:projb}.}
  \label{fig:projf}
\end{figure}

\subsection{Kobe strategy matrix}
Figure~\ref{fig:k2sm} is $\PG$ by catch and year for the same six series.
The two MVLN panels are the \nDraw{}-draw sample described in
Section~\ref{sec:mvln}. SS3-MCMC, FLasher and both JABBA3 panels are
proportions of their own draws. Table~\ref{tab:tacfour} uses the four
projections only.

\begin{figure}[H]
  \centering
  \incfig[width=\textwidth]{k2sm-1.png}
  \caption{Kobe strategy matrix: $\PG$ (\%) by catch and year. recDev2 and
    FLasher are the two extra series from the shared knit.}
  \label{fig:k2sm}
\end{figure}

\subsection{Risk-equivalent catch}
At zero catch, $\PG$ in 2045 was \greenZeroMvlnFortyFive\% (SS3-MVLN),
\greenZeroMcmcFortyFive\% (SS3-MCMC), \greenZeroJabbaFortyFive\% (JABBA3)
and \greenZeroJabbaDynFortyFive\% (JABBA3 dyn). In 2070 the same catches
give \greenZeroMvlnSeventy\%, \greenZeroMcmcSeventy\%,
\greenZeroJabbaSeventy\% and \greenZeroJabbaDynSeventy\%.

The catch at $\PG = 0.6$ in 2070 is \tacMvlnSeventySixty~t under SS3-MVLN
and \tacMcmcSeventySixty~t under SS3-MCMC. At $\PG = 0.7$ in 2070 the
catches are \tacMvlnSeventySeventy~t and \tacMcmcSeventySeventy~t
(Table~\ref{tab:tacfour}, Figure~\ref{fig:risk}). Both JABBA3 projections
stay above 0.6 and 0.7 at every catch on the grid, including 2{,}300~t.
At 2{,}300~t in 2070, $\PG$ is 72\% with constant process error and 80\%
with state-dependent process error. The risk-equivalent catch for both
JABBA3 projections is therefore above 2{,}300~t, and the table records the
top of the grid. The SS3-MVLN catch at 60\% in 2070 is at least
\tacRangeSixtySeventy~t below that ceiling. State-dependent process error
does not pull the surplus-production tier toward the SS3 catches. It is the
more optimistic of the two JABBA3 projections.

The SS3-MVLN and SS3-MCMC columns should not be read as a pure method
contrast (Section~4). The JABBA3 columns are the structural contrast: the
same nominal $\PG$ does not select the same catch once the projection drops
the age-structured lag.

\begin{figure}[H]
  \centering
  \incfig[width=\textwidth]{risk-fig-1.png}
  \caption{$\PG$ against catch in 2030, 2045 and 2070. Dashed lines are
    $\PG = 0.6$ and $0.7$. The catch at which a curve crosses a line is the
    risk-equivalent catch. recDev2 and FLasher are drawn with the four
    projections; Table~\ref{tab:tacfour} reports the four only.}
  \label{fig:risk}
\end{figure}

\begin{table}[htbp]
\caption{\label{tab:tacfour}Catch (t) at which $\PG$ equals the target.
A value of 2{,}300~t is the top of the grid: the target is still met there,
so the crossing lies above the grid.}
\centering
\begin{tabular}[t]{lrrrr}
\toprule
Projection & 2045 @ 60\% & 2070 @ 60\% & 2045 @ 70\% & 2070 @ 70\%\\
\midrule
SS3-MVLN & \tacMvlnFortyFiveSixty & \tacMvlnSeventySixty &
  \tacMvlnFortyFiveSeventy & \tacMvlnSeventySeventy\\
SS3-MCMC & \tacMcmcFortyFiveSixty & \tacMcmcSeventySixty &
  \tacMcmcFortyFiveSeventy & \tacMcmcSeventySeventy\\
JABBA3 & \tacJabbaFortyFiveSixty & \tacJabbaSeventySixty &
  \tacJabbaFortyFiveSeventy & \tacJabbaSeventySeventy\\
JABBA3 dyn & \tacJabbaDynFortyFiveSixty & \tacJabbaDynSeventySixty &
  \tacJabbaDynFortyFiveSeventy & \tacJabbaDynSeventySeventy\\
\bottomrule
\end{tabular}
\end{table}

\subsection{Rebuilding year}
Figure~\ref{fig:rebuild} plots the rebuilding year against catch for each
target. Vertical separation at a given catch is the difference in rebuilding
time. Both JABBA3 curves lie on the first projection year across this grid,
because $\PG$ is already above the target, so the catch inversion in
Table~\ref{tab:tacfour} is the summary for those two projections. The
recDev2 MVLN and FLasher curves are drawn with the four projections and are
not part of that table.

\begin{figure}[H]
  \centering
  \incfig[width=\textwidth]{rebuild-fig-1.png}
  \caption{Rebuilding year against catch. The year is the last crossing of
    the target $\PG$, given that the series ends at or above the target.
    Line type is the target (0.6 or 0.7). A curve along 2025 is a series
    that already meets the target at the start of the projection.}
  \label{fig:rebuild}
\end{figure}

\section{Discussion}

The four projections answer one question. At a fixed $\PG$ and year, what
catch does each projection return, and at a fixed catch, in what year is
the target met. They do not replace the SS3 projection endorsed by the
Shark WG \citep{Liniers2026}.

On this grid the SS3-MVLN catch at 60\% in 2070 is \tacMvlnSeventySixty~t.
Both JABBA3 projections are still above 60\% and 70\% at 2{,}300~t. The gap
is structural. JABBA3 has no age-structured lag between exploitable biomass
and spawners \citep{Winker2020SMA}, and state-dependent process error makes
that tier more optimistic, not closer to SS3. The bridge is included so
that disagreement is visible.

The SS3-MCMC catches are lower than the SS3-MVLN catches. The MCMC was run
with \texttt{do\_recdev} = 2 and the MVLN base case with
\texttt{do\_recdev} = 1, and the reference points differ by about a third
(Table~\ref{tab:refpts}). That contrast is not a test of the delta method
against MCMC.

The MVLN probability has the closed form in Section~\ref{sec:mvln}. Nothing
in that expression has to be estimated: the means, standard deviations and
correlation are already in the SS3 covariance. The 2026 matrix was
nevertheless filled by drawing \linN{} times from that distribution, up
from 15{,}000 ``to stabilize results'' \citep{Liniers2026}. That simulation
was unnecessary. Drawing from a distribution whose probability is already
known does not calculate the probability. It replaces an exact value with a
sample mean, and then more draws are used to average the noise back out.
Appendix~\ref{app:mvln} shows that the closed form reproduces the published
cells, and that the differences which remain are that noise. The catches in
this paper are themselves a sample of \nDraw{} draws, reported as such.
The matrix should be evaluated, not simulated.

Two further series are drawn in the figures. The recDev2 MVLN uses the
recruitment constraint of the MCMC model, so it is the partner for a later
comparison that is not confounded with \texttt{do\_recdev}. FLasher carries
block-resampled recruitment residuals on an MCMC terminal state. Neither
series is in Table~\ref{tab:tacfour}.

\section{Conclusions}
\begin{enumerate}[leftmargin=*]
  \item Use one catch grid, one pair of target probabilities and one
    definition of green for SS3-MVLN, SS3-MCMC, JABBA3 and JABBA3 with
    state-dependent process error.
  \item Fill an MVLN Kobe matrix in closed form. Do not derive those
    probabilities by Monte Carlo. The 2026 matrix was filled by simulation
    \citep{Liniers2026}; that step was unnecessary, because $\PG$ was
    already the bivariate normal probability of the reported covariance.
    Fill SS3-MCMC and both JABBA3 cells by counting draws, because those
    distributions have no such form.
  \item Report the catch at $\PG = 0.6$ and $0.7$ in 2045 and 2070, and the
    rebuilding year. On this SMA-N grid the JABBA3 catches lie above
    2{,}300~t; the SS3-MVLN catch at 60\% in 2070 is \tacMvlnSeventySixty~t.
  \item Do not treat the SS3-MVLN versus SS3-MCMC difference as a method
    test until both use the same \texttt{do\_recdev}. Do not treat the
    JABBA3 bridge as Shark WG catch advice.
\end{enumerate}

\appendix
\renewcommand\thesection{Appendix \Alph{section}.}
\section{Closed-form MVLN against the published Kobe matrix}
\label{app:mvln}

The 2026 SMA-N Kobe matrix was filled by Monte Carlo
\citep{Liniers2026}. \linN{} draws were taken from the multivariate
log-normal, increased from 15{,}000 ``to stabilize results''. Under the
model in Section~\ref{sec:mvln} there was nothing to stabilize. The means,
standard deviations and correlation are the SS3 covariance, and $\PG$ is
the bivariate normal probability of those five numbers. A draw from that
distribution does not estimate $\PG$. It substitutes a binomial proportion
for a probability that is already known, and the only effect of a larger
sample is to hide the substitution. A probability near one half, estimated
from \linN{} draws, still has a standard error of about 0.35 percentage
points, enough to move a rounded cell across a boundary. The simulation was
unnecessary. It added no information, and the discrepancies below are the
noise it introduced.

To confirm that the closed form reproduces the published table, \linCells{}
cells of that matrix were recomputed from the same SS3 runs
(Table~\ref{tab:liniers}). Matching the published cells requires the
conventions in \texttt{ss3diags2::ss3mvln}, as in
\texttt{ss3diags::SSdeltaMVLN}: end-of-year SSF, which is the SS3
\texttt{Bratio} of the following year paired with $F$ of the current year,
and the log-scale correlation set equal to the natural-scale correlation.
With those conventions the analytical values lie within one percentage point
of the published values in \linWithinOne{} of the \linCells{} cells, and
within \linMaxRounded{} percentage points in every cell. The remainder is
rounding plus the Monte Carlo error of the \linN{} draws.

The projections in the body of this paper do not use that natural-scale
correlation. They use $\log(1 + \rho\,\mathrm{CV}_B\mathrm{CV}_F)$. In these
cells the two conventions change $\PG$ by at most \trDiff{} percentage
points. Both are analytical. The natural-scale correlation is used only so
that Table~\ref{tab:liniers} is a comparison with the published matrix, not
a second projection.

\begin{table}[H]

\caption{\label{tab:liniers}Analytical MVLN Kobe probabilities (\textbackslash{}\%) from the SS3 base case, with the 20,000-draw values of Liniers et al. (2026, Table 1) in brackets.}
\centering
\begin{tabular}[t]{lrllll}
\toprule
Quantity & Catch (t) & 2030 & 2040 & 2050 & 2070\\
\midrule
P(F<=Fmsy) & 0 & 100 (100) & 100 (100) & 100 (100) & 100 (100)\\
P(F<=Fmsy) & 250 & 100 (100) & 100 (100) & 100 (100) & 100 (100)\\
P(F<=Fmsy) & 1200 & 98 (98) & 99 (98) & 99 (99) & 100 (100)\\
P(F<=Fmsy) & 1883 & 44 (44) & 41 (40) & 42 (42) & 43 (43)\\
P(F<=Fmsy) & 2300 & 13 (13) & 11 (11) & 11 (11) & 12 (12)\\
\addlinespace
P(SSF>=SSFmsy) & 0 & 47 (45) & 82 (86) & 99 (99) & 100 (100)\\
P(SSF>=SSFmsy) & 250 & 46 (44) & 77 (80) & 97 (97) & 100 (100)\\
P(SSF>=SSFmsy) & 1200 & 43 (40) & 52 (54) & 68 (69) & 87 (87)\\
P(SSF>=SSFmsy) & 1883 & 41 (37) & 35 (35) & 37 (37) & 40 (40)\\
P(SSF>=SSFmsy) & 2300 & 39 (35) & 26 (26) & 23 (23) & 20 (20)\\
\addlinespace
P(green) & 0 & 47 (45) & 82 (86) & 99 (99) & 100 (100)\\
P(green) & 250 & 46 (44) & 77 (80) & 97 (97) & 100 (100)\\
P(green) & 1200 & 43 (40) & 52 (54) & 68 (69) & 87 (87)\\
P(green) & 1883 & 33 (31) & 29 (30) & 32 (33) & 35 (36)\\
P(green) & 2300 & 12 (12) & 10 (10) & 10 (11) & 11 (12)\\
\bottomrule
\end{tabular}
\end{table}


\bibliographystyle{plainnat}
\bibliography{refs}

\end{document}
